\section{An all-order renewal-root cut from the positive primitive}
This is an auxiliary closure theorem, not a proof of general actual-arm
C or original one-descent coefficient unimodality. All renewal orders
and delay parameters in the theorem are unbounded.

Fix $t\ge2$, let $0<q<1$ be the unique solution of
$q+\cdots+q^t=1$, and put $\lambda=q^{-1}$ and $p=t-1$. Set
\[
 Q_t=1-z-\cdots-z^t=(1-\lambda z)B_t(z),\qquad
 (B_t)_i=\sum_{k=1}^{t-i}q^k\quad(0\le i\le t-1).
\]
In particular $(B_t)_0=1$. All sequences below have zero extension
outside their specified finite intervals. A sequence is LC when its
adjacent $2$ by $2$ minors are nonnegative and its positive support is
an interval; the identically zero sequence is permitted.

\begin{theorem}[All-order cut]
Let $J$ be finite, nonnegative, unimodal, and positive on an interval
starting at zero. Let $M$ be its first mode. Assume
\[
 e_n=(J_n-J_{n-1})^+\boldsymbol1_{n\ge1}
\]
is LC. For any $\ell\ge0$ define, on integer indices,
\[
 x_n=\sum_{j\ge1}q^jJ_{n+j},\quad Y=[p]x,\quad
 A_n=\bigl(Y_n-(z^\ell B_tJ)_n\bigr)^+.
\]
Then $A$ is LC at every centre $\ell+t\le n<M$.
If $E=B_tJ+R$ with $\deg R\le t-2$, the same conclusion holds
with $z^\ell E$ in place of $z^\ell B_tJ$, on exactly that band.
No claim beyond the first mode of $J$ is made.
\end{theorem}

\begin{lemma}[A complete kernel calculation]
Let $X_j=q^{-j}$ for $j\le-1$ and zero otherwise. The bilateral
kernel
\[
 K=\frac{[p]X-z^\ell B_t}{1-z}
\]
has positive support $(-\infty,U]$, $U=\ell+t-2$, and is LC.
Here division by $1-z$ means cumulative summation from the negative
infinite end; the sum converges. Its negative head is exactly
\[
 K_j=\frac{[p](q)}{1-q}q^{-j}\quad(j\le-1).
\]
Writing $L_n=K_{U-n}$ for $n\ge0$ gives the complete ordinary
generating function
\[
 L(z)=\frac{q\{[p](z)[h](z)+z^p\}}{1-qz},\qquad h=\ell+2.
 \tag{1}
\]
\end{lemma}
\begin{proof}
Cumulative summation gives explicitly
\[
 (1-q)K_j=\sum_{s=0}^{p-1}q^{\max(1,s-j)}
 -(1-q)\sum_{i=0}^{\min(t-1,j-\ell)}(B_t)_i,
 \tag{2}
\]
where the last sum is empty when $j<\ell$. Equivalently, reversing
the identity
$(1-\lambda z)(1-z)K=-[p]-z^\ell Q_t$ gives
\[
 \lambda(1-qz)(1-z)L=[p+1]-z^{p+1}-z^h[p]
       =(1-z)([p][h]+z^p).
\]
These are convergent bilateral identities on $|z|>q$ or, equivalently,
finite coefficient identities from (2); no formal infinite Laurent
series multiplication is required. This proves (1), its support, and
its negative head.

It remains to prove LC of $F=g/(1-qz)$ for $g=[p][h]+z^p$.
The root equation implies $1/2\le q\le2/3$: its finite geometric
identity gives $1-q=q-q^{t+1}<q$, while $q+q^2\le1$ gives the
upper bound. We consider both possible width orders.

If $2\le h\le p$, the coefficients of $g$ are the initial ramp
$1,2,\ldots,h$, a plateau at $h$ extending through index $p$,
and then $h-2,h-3,\ldots,1$. For $h=2,3,4$ this is LC:
the only potentially exceptional junction is
$(h-2)^2\ge h\max(0,h-3)$. Geometric convolution preserves LC.
For $h\ge5$, compare the sequence through index $p$ with the
globally LC source $\min(n+1,h)$ and put
\[
 V=\frac h{1-q},\quad
 R=\frac{q^{p-h+2}(1-q^h)}{(1-q)^2}\ge0.
\]
Then
\[
 F_p=V-R,\quad F_{p+1}=V-2-qR,\quad
 F_{p+2}=V-3-2q-q^2R.
\]
The centre-$p$ minor improves when the next source is reduced by
two. The only new centre satisfies the exact identity
\[
 F_{p+1}^2-F_pF_{p+2}
 =4+V(2q-1)+R\{(h-2)(1-q)-1\}\ge4.
 \tag{3}
\]
The last sign uses $h\ge5$ and $q\le2/3$. Thereafter $g$ is a
decreasing linear ramp and is LC. The following exact propagation
identity pays every remaining positive-source centre. If
$B=qA+g$ and $C=qB+h$, then
\[
 g(Ch-Bk)=qh(Bg-Ah)+B(h^2-gk).                 \tag{4}
\]
Apply (4) to consecutive source coefficients $g,h,k$, using the
previous minor and the next source minor. After the source ends
the sequence is geometric; its joining minor is nonnegative and
its later minors vanish.

If $h>p$, the source through $p$ is the complete ramp $n+1$.
The root equation implies
\[
 (1-q)F_p=p.
 \tag{5}
\]
Indeed $(1-q)\sum_{i=0}^pq^{p-i}(i+1)
=p+1-\sum_{j=1}^{p+1}q^j=p$. Thus $F$ reaches
$V=p/(1-q)$ exactly at $p$ and is constant through $h-1$.
For $p\ge2$, the first decreasing values are $V-1,V-2-q$,
whose new minor is $1+qV>0$. The ramp-to-plateau joins are LC;
(4) pays the remaining decreasing-source tail. For $p=1$ the
source ends at $h-1$ and the following tail is already geometric.
This covers every $p\ge1,h\ge2$ and proves the lemma.
\end{proof}

\begin{proof}[Proof of the theorem]
Put $d=(1-z)J$, including its final negative entry. The finite
convolution identity is
\[
                   Y-z^\ell B_tJ=K*d.                    \tag{6}
\]
Every negative $d_k$ has $k>M$ by unimodality. For $n\le M$
these terms use only the geometric negative head of $K$, giving
exactly $-T\lambda^n$ for one constant $T\ge0$. If $n>U$,
the $d_0$ term is zero as well. Hence on $U<n\le M$,
\[
             (K*d)_n=(K*e)_n-T\lambda^n.                \tag{7}
\]
The first term is LC and nonnegative: truncate the negative
kernel tail below the finite read band and use ordinary LC
convolution. The truncations agree on every needed finite band.
After division by $\lambda^n$, subtract $T$ and take positive
parts. This preserves LC, including all zero boundaries, by the
registered constant-cut lemma in \texttt{GoldenPrimitiveCut.lean}.
Multiplication back by $\lambda^n$ preserves LC. All three reads
at $\ell+t\le n<M$ lie in $U<n-1\le n+1\le M$.
Finally, at these centres all three source indices are at least
$t-1$, so a remainder of degree at most $t-2$ has no contribution.
\end{proof}

\paragraph{Checks and exact scope.}
\texttt{RenewalRootPrimitiveCut.lean} checks nine standard-axiom
reports: the notch identity and its complete large-width payment;
all small widths; the unbounded ramp fixed-point induction; all
renewal root bounds; the full interval-product factorization;
the exact three-read band; tail-minor propagation; and the first
decreasing-ramp minor. Kernel reversal, interval-support convolution,
and signed-tail assembly are written arguments, not an end-to-end
formalization of composition hosts. Exact root-ring tests are QA
only and introduce no finite premise into the theorem.

For $t=2$ this recovers the golden cut with $Y=x$; for $t=3$
it gives $Y=x+zx$ and the proper quadratic tribonacci denominator.
The factor $[t-1]$ is required. The theorem does not identify an
actual binary/ternary host with this model and does not discard
the actual extra source polynomial.
